\documentclass{article}
\usepackage{graphicx} % Required for inserting images

\title{Lesson 10 Homework - Calculus}
\author{William Wang}
\date{September 2026}

\begin{document}

\maketitle

\section{}
We have the function 
\[f(x) = \sqrt{|4-x|}\]
To see whether it is differentiable at \(4\) we want to see if the right hand derivative is equal to the left hand derivative. Right hand derivative
\[\lim_{h\to 0^+} \frac{f(4+h) - f(4)}{h} = \lim_{h\to 0^+} \frac{\sqrt{|4-4-h|} - \sqrt{|4|}}{h} = \lim_{h\to 0^+}\frac{\sqrt{h}}{h}\]
\[ = \lim_{h\to 0^+} \frac{1}{\sqrt{h}} = \infty\]
Then the left hand derivative is 
\[\lim_{h\to 0^-} \frac{f(4+h) - f(4)}{h} = \lim_{h\to 0^-} \frac{\sqrt{|h|}}{h}\]
since \(h < 0\), 
\[= \lim_{h\to 0^-} \frac{\sqrt{-h}}{h}\]
since \(h < 0\), we can write \(h\) as \(-(-h)\). Then 
\[\lim_{h\to 0^-} \frac{\sqrt{-h}}{h} = \lim_{h\to 0} \frac{\sqrt{-h}}{- (\sqrt{-h})^2} = \lim_{h\to 0} \frac{1}{-\sqrt{-h}} = -\infty\]
Graphically, as $x$ approaches $4$ from the right, the tangent slopes become infinitely steep and positive. As $x$ approaches $4$ from the left, the tangent slopes become infinitely steep and negative. This forms a shape that looks like 2 semi circles tangent to each other and opening away from each other. 
\end{document}
